\chapter{Alineamiento con los Objetivos de Desarrollo Sostenible (ODS)}
\label{sec:obj_ds}

El presente Trabajo de Fin de Grado no solo aborda un reto técnico complejo en el ámbito de la Ingeniería de Datos y el Machine Learning, sino que nace con la vocación de aportar valor real a la sociedad, alineándose de forma directa con la Agenda 2030 de las Naciones Unidas. En concreto, el desarrollo de este sistema de ingesta, imputación y predicción de ocupación de plazas de aparcamiento impacta directamente en los siguientes Objetivos de Desarrollo Sostenible:

En este nuevo paradigma, los sensores IoT actúan como el ``sistema nervioso'' de la ciudad. Estos dispositivos, embebidos en el asfalto, semáforos, farolas o contenedores, son capaces de capturar, procesar y transmitir parámetros físicos en tiempo real. La verdadera disrupción tecnológica del IoT urbano no reside únicamente en la recolección aislada de datos, sino en la capacidad de generar un flujo continuo de información a gran escala (Big Data). Esto permite a las administraciones públicas y a los sistemas automatizados pasar de una gobernanza reactiva (actuar cuando hay un problema) a una proactiva y predictiva.

\subsection{ODS 11: Ciudades y Comunidades Sostenibles}
\label{sec:obj_ds}
El objetivo principal de este proyecto es optimizar la movilidad urbana mediante el uso inteligente de los datos. En las grandes ciudades, una parte significativa del tráfico rodado está compuesto por el denominado ``tráfico de agitación'': vehículos circulando sin rumbo fijo en busca de aparcamiento. Consecuentemente, al predecir con exactitud la disponibilidad de plazas usando redes neuronales sobre datos de sensores IoT, este sistema proporciona la base tecnológica necesaria para informar a los conductores en tiempo real y optimizar las rutas. Esto fomenta ciudades más habitables, reduce la congestión de las vías públicas y mejora la planificación del espacio urbano por parte de las administraciones públicas.

\subsection{ODS 9: Industria, Innovación e Infraestructura}
\label{sec:obj_ds}
La modernización de las infraestructuras urbanas requiere sistemas capaces de lidiar con la imperfección del mundo físico. Los sensores de las calles fallan, se quedan sin batería o pierden conexión, generando vacíos de información. En este proyecto, se aporta una solución innovadora al diseñar una arquitectura de datos robusta combinada con modelos de Deep Learning de última generación (PyPOTS). La capacidad del modelo para aprender patrones espacio-temporales e imputar valores faltantes demuestra cómo es posible hacer que las infraestructuras de hardware existentes de sensores IoT sean mucho más fiables, resilientes y escalables sin necesidad de incurrir en costes masivos de mantenimiento.

\subsection{ODS 13: Acción por el Clima}
\label{sec:obj_ds}
Por último, cabe destacar el impacto en contaminación por el sector del transporte, uno de los mayores emisores de gases de efecto invernadero a nivel global. Cualquier ineficiencia en la conducción se traduce directamente en contaminación evitable. De forma indirecta pero medible, facilitar que un vehículo se dirija directamente a una zona con alta probabilidad de aparcamiento disponible reduce drásticamente el tiempo de motor en marcha. Esta reducción del tiempo de búsqueda minimiza las emisiones de CO₂, NOx y partículas en suspensión en los núcleos urbanos, contribuyendo de manera efectiva a la mitigación del cambio climático y mejorando la calidad del aire que respiran los ciudadanos.